\chapter*{Résumé}
\addcontentsline{toc}{chapter}{Résumé}

Cette thèse traite de la prédiction faiblement supervisée de
mutations à partir d'images H\&E (colorées à
l'hématoxyline-éosine) gigapixel d'adénocarcinome pulmonaire
(LUAD, le sous-type le plus fréquent du cancer du poumon) :
une seule étiquette binaire par gène est attribuée à environ
$10^9$ pixels, sans annotation au niveau des sous-régions.
Les systèmes existants requièrent des dizaines de milliers de
coupes propriétaires et n'offrent aucune explication du
\emph{pourquoi} d'une mutation prédite --- un obstacle critique à
la confiance clinique. Le problème sous-jacent est lui-même
\emph{stochastique, plusieurs-à-plusieurs et faiblement étiqueté} :
une même morphologie peut correspondre à plusieurs génotypes et
inversement, de sorte que la compétence d'un modèle doit être
argumentée gène par gène et coupe par coupe, et non à partir
d'une seule métrique de cohorte.

L'hypothèse centrale est qu'une \emph{connaissance structurée du
domaine} --- en l'occurrence, les six motifs de croissance
histologiques du LUAD définis par les experts
(lepidic, acinar, papillary, micropapillary, solid, cribriform),
encodés via la théorie des ensembles flous --- peut être injectée
dans un pipeline d'apprentissage profond afin d'atteindre à la
fois une prédiction de mutations compétitive en condition de
pénurie de données ($<\!700$ coupes publiques) et de rendre les
prédictions résultantes \emph{interprétables} grâce à un cadre
d'inspection multi-niveaux cliniquement vérifiable --- et, à
travers ce cadre, \emph{explicables} par des justifications en
langage naturel ancrées dans un graphe de connaissances --- à un
degré qu'aucune approche boîte noire de bout en bout n'offre ;
le verdict est \emph{conditionnel au gène}.
Suivant la démarche de la Design Science Research, trois
artefacts et un système intégrateur (\lchai{},
\emph{Lung Cancer Histologic Analysis with AI}) sont conçus et
évalués.

\textbf{Artefact~1 (FuzzyArcLoss~V2)} est une modulation floue
échantillon-par-échantillon de la marge angulaire d'ArcFace,
respectant les transitions morphologiques graduelles entre
classes de motifs ; il se classe premier parmi 18~fonctions de
perte (macro-F1 $= 92{,}31\,\%\!\pm\!2{,}04$, $p\!=\!0{,}001$,
$d$ de Cohen $\!=\!1{,}00$). Son avantage est
\emph{dépendant des frontières morphologiques, et non de la
taille des classes} : le gain se matérialise sur les classes
voisines visuellement ambiguës, indépendamment de leur taille
--- une règle de transfert qui se généralise au grade de
Gleason, à la cytologie cervico-utérine et à la dermatopathologie.

\textbf{Artefact~2 (PI-ABMIL)} concatène un \emph{embedding}
CTransPath de dimension~512 avec les six probabilités floues de
motifs issues de l'Artefact~1, avant agrégation par
attention~\emph{gated}, ancrant le raisonnement au niveau de la
coupe dans des motifs cliniquement validés.

\textbf{Artefact~3 (FC-MIL)} remplace l'attention additive par
une intégrale de Choquet 2-additive apprise, dont les
21~paramètres (six valeurs de Shapley, quinze interactions par
paires) capturent les interactions supra-additives entre
catégories morphologiques que l'attention standard ne peut
représenter.

L'évaluation sur TCGA-LUAD (\emph{The Cancer Genome Atlas Lung
Adenocarcinoma} ; 687~coupes, 668~patients, validation croisée
5-fold stratifiée par patient) sur six mutations conductrices
(\textit{TP53}, \textit{EGFR}, \textit{KRAS}, \textit{STK11},
\textit{KEAP1}, \textit{RBM10} ; prévalence décroissante)
montre que l'utilité de la connaissance des motifs est
\emph{conditionnelle au gène et à l'agrégateur}, et se décompose
en trois régimes structurels :

\textbf{Régime~(i) --- saturation de l'encodeur.}
L'encodeur pré-entraîné capture déjà la morphologie pertinente
et le canal des motifs n'apporte rien (\textit{TP53},
\textit{EGFR}).

\textbf{Régime~(ii) --- co-occurrence inter-motifs.}
Le signal discriminant réside dans les interactions par paires
de motifs, et non dans un motif unique ; seul un agrégateur
lisant les activations conjointes le récupère (\textit{KRAS},
via l'intégrale de Choquet de FC-MIL).

\textbf{Régime~(iii) --- hors de l'ontologie.}
Aucune représentation à l'intérieur du simplexe à six motifs
(ANORAK~\cite{anorak2024}) ne porte le signal ; le régime se
décompose en trois mécanismes nommés --- \emph{label-space
mismatch}~\cite{lipton2018labelshift}, \emph{no morphological
projection} et
\emph{collapse-to-prior}~\cite{bowman2016posteriorcollapse} ---,
qui ensemble bornent ce que la prédiction de mutations à partir
de H\&E seule peut atteindre (\textit{STK11}, \textit{KEAP1},
\textit{RBM10}).

Pour rendre ces limites visibles, \lchai{} instrumente une pile
d'interprétabilité à six niveaux (attention, surimpressions de
motifs, SHAP, valeurs de Shapley issues de Choquet, ablation de
canaux, étiquettes linguistiques floues) et un agent
d'explication LLM ancré dans un graphe de connaissances curé,
construit à partir des ontologies NCIt et MONDO : les étiquettes
floues contraignent \emph{comment} le LLM s'exprime, le graphe
contraint \emph{ce qu'}il peut affirmer --- un mécanisme dual
contre l'hallucination.
Un drapeau complémentaire au niveau cohorte
\emph{Conclusive\,/\,Inconclusive} (dérivé de l'AUROC) régit la
voie d'explication par coupe, séparant la présence d'un signal au
niveau de la population de la confiance morphologique au niveau
de chaque coupe.

L'hypothèse est \emph{partiellement confirmée} : la connaissance
des motifs apporte une interprétabilité qu'aucun système
antérieur de prédiction de mutations n'offre, ainsi qu'un
véritable gain prédictif inter-motifs sur \textit{KRAS} ; pour
les autres gènes, la borne réside dans la \emph{représentation
d'entrée} à six motifs elle-même (l'espace d'étiquettes ANORAK)
et non dans le modèle aval
(perte\,/\,fusion de caractéristiques\,/\,agrégateur).
Les revendications principales de cette thèse ne sont donc pas
une précision de qualité clinique --- des cohortes à l'échelle
des modèles de fondation, comportant plus de 60\,000 coupes,
existent déjà à cette fin --- mais les deux suivantes,
distillées des six contributions fines C1--C6 listées au
chapitre~\ref{ch:introduction} :

\textbf{Revendication~(i) --- méthodologique.}
La logique floue, appliquée de manière cohérente à la perte, à
la représentation et à l'agrégation, rend transparent un
pipeline reposant sur des données publiques, à travers une pile
d'interprétabilité à six niveaux ancrée dans un graphe de
connaissances.

\textbf{Revendication~(ii) --- empirique/analytique.}
La taxonomie gène par gène issue des ablations par coupe sur
TCGA-LUAD borne ce qu'un modèle reposant uniquement sur la
morphologie peut atteindre à partir de H\&E seule, et identifie
les gènes pour lesquels la fusion de sources de données
supplémentaires (par exemple profils d'expression génique ou
radiologie), plutôt qu'un raffinement architectural
supplémentaire, constitue l'étape naturelle suivante.

Ensemble, ces revendications positionnent l'approche comme un
paradigme complémentaire à la force brute des modèles de
fondation --- un paradigme dans lequel une pathologiste ou un
pathologiste peut comprendre, vérifier et \emph{corriger} le
raisonnement du modèle, et dans lequel l'interprétabilité est la
plus précieuse précisément sur les cas que le prédicteur prédit
de manière erronée.

\paragraph{Mots-clés.}
Logique floue, perte à marge angulaire, apprentissage à
instances multiples, intégrale de Choquet, valeurs de Shapley,
motifs histopathologiques, adénocarcinome pulmonaire, prédiction
de mutations, fossé génotype--phénotype, Design Science Research,
pathologie computationnelle, IA explicable.
